\section{What TOLL-PIBT changes, and why it is different}
\label{sec:contribution}

\begin{lede}
Four changes, three of them consequences of one idea: a high-level traffic
rule may reorder what a local planner considers, and must never remove it.
Each is stated as an inequality or a property, with the measurement that tests
it.
\end{lede}

The architecture is two layers, and the division between them is the whole
design (\Cref{sec:architecture}). The upper layer decides what each robot
\emph{wants}: which task, which waypoint, which route, which direction its
aisle should be flowing, who plans first. The lower layer --- PIBT, unmodified
--- decides what is \emph{allowed}. The upper layer may rank the lower layer's
options. It may not delete them. Everything below follows from taking that
sentence literally.

\subsection{Direction as a term in the objective, not a constraint on the domain}
\label{sec:contribution-soft}

\subsubsection{The two implementations}

Let $C_i = N(\pos{i}{t}) \cup \{\pos{i}{t}\}$ be robot $i$'s candidate set and
let $\mathrm{safe}_i(v)$ hold when moving to $v$ creates neither conflict of
\Cref{eq:vertex-conflict,eq:swap-conflict}. Write
$\mathrm{opp}_i(v)$ for the predicate ``moving to $v$ opposes the committed
direction of the aisle being entered'' (\Cref{sec:aisle-legality}). The two
ways to enforce a one-way rule are:
\begin{align}
  F^{\mathrm{hard}}_i &= \{\, v \in C_i \;:\; \mathrm{safe}_i(v) \wedge \neg\,\mathrm{opp}_i(v) \,\},
  \label{eq:hard}\\
  F^{\mathrm{soft}}_i &= \{\, v \in C_i \;:\; \mathrm{safe}_i(v) \,\},
  \qquad
  \score \;\mathrel{-}=\; \zeta \cdot \ind{\mathrm{opp}_i(v)}.
  \label{eq:soft}
\end{align}
\Cref{eq:hard} is what the original specification for this system called for,
and what a one-way rule ordinarily means. \Cref{eq:soft} is what TOLL-PIBT
does: the move stays in the set, sorted to the back, and is taken when nothing
better is available.

\subsubsection{Why the difference is not cosmetic}

Under a path planner the two are nearly equivalent --- a forbidden edge simply
raises the cost of the route through it, and the planner routes around. Under
PIBT they are not, because PIBT's mechanism for resolving a blockage is to
\emph{push}, and pushing needs somewhere to push to.

\begin{proposition}[Candidate preservation]
\label{prop:preservation}
For every robot $i$ and timestep $t$, $F^{\mathrm{hard}}_i \subseteq
F^{\mathrm{soft}}_i$, and $F^{\mathrm{soft}}_i$ is exactly the set PIBT would
have without any aisle layer at all. Consequently every joint move reachable
under \Cref{eq:hard} is reachable under \Cref{eq:soft}, and the aisle layer
under \Cref{eq:soft} cannot reduce the set of configurations the planner can
escape from.
\end{proposition}

The converse fails, and it fails exactly where it hurts. Consider a
single-file aisle $a = (u_1, \dots, u_m)$ committed \textsc{forward}
(increasing index), occupied by robots $i_1, \dots, i_m$, one per cell, with
$u_m$'s forward neighbour occupied by a robot that cannot move this timestep.
Under \Cref{eq:hard}, each $i_k$ has exactly two candidates: $u_{k+1}$, which
is occupied, and staying put; the reverse move $u_{k-1}$ has been deleted.
Priority inheritance therefore walks the chain $i_1 \to i_2 \to \dots \to i_m$,
finds no robot with a free legal move, backtracks all the way, and every robot
stays. The configuration is unchanged next timestep, and the timestep after
that. Under \Cref{eq:soft}, $i_1$ retains $u_0$ at a cost of $\zeta$, the chain
terminates by pushing one robot backwards out of the aisle, and the aisle
drains.

This is not a hypothetical. Instrumented on \variant{warehouse\_corridors}
under \Cref{eq:hard}: all 35 robots stalled by $t = 120$, four aisles locked
\textsc{reverse}, the identical global state still present at $t = 199$. The
second comparison animation in \code{docs/gifs/} is this pair of runs side by
side.

\subsubsection{Sizing \texorpdfstring{$\zeta$}{zeta}}
\label{sec:contribution-zeta}

Pricing only works if the price is right. The candidate score
(\Cref{sec:score-full}) is
\begin{multline}
  \score = \underbrace{\alpha \cdot \mathrm{progress}}_{\alpha = 10}
    + \underbrace{\beta_i \cdot \ind{\text{preferred direction}}}_{\beta \leq 3}
    + \underbrace{\gamma_i \cdot \ind{\text{stays in aisle}}}_{\gamma \leq 2}
    - \underbrace{\lambda_{\mathrm{turn}} P_{\mathrm{turn}}}_{\leq 1}\\
    - \underbrace{\mu C_i(v)}_{\leq 1}
    - \underbrace{\nu P_{\mathrm{wait}}}_{\leq 0.2}
    - \underbrace{\xi \ind{\text{bottleneck}}}_{\leq 1}
    - \underbrace{\zeta \ind{\mathrm{opp}_i(v)}}_{\zeta = 8},
\end{multline}
and $\zeta$ is placed by two inequalities.

\paragraph{$\zeta < \alpha$: progress can still buy its way through.}
A counterflow move that closes one step of distance scores
$\alpha - \zeta = +2$ before minor terms; staying in place scores
$-\nu = -0.2$. So a robot in a corridor whose only progressing move is against
the flow \emph{takes it}, and pays. This is the inequality that makes
\Cref{prop:preservation} operational rather than merely set-theoretic: the
retained candidate is not just present, it is reachable.

\paragraph{$\zeta$ above every other preference: direction still governs.}
Individually, $\zeta = 8$ exceeds the largest competing soft term
($\beta_{\max} = 3$) by $2.7\times$, and exceeds the sum of both positive
preference terms ($\beta + \gamma = 5$). It does not exceed the arithmetic sum
of all six non-progress terms at their simultaneous extremes ($8.2$), and we
state that plainly rather than claiming a domination that the numbers do not
support: the guarantee is that no ordinary combination of preferences reverses
a direction decision, not that no contrived one can. What is guaranteed
exactly is the ordering against progress, above.

\subsubsection{What it is worth}

Single-factor comparisons (\code{config.PAIRED\_DESIGNS}), same seeds, same
maps, the same direction decisions, differing only in \Cref{eq:hard} versus
\Cref{eq:soft}: $1.9\times$ throughput on \variant{warehouse\_corridors},
$2.5\times$ on \variant{warehouse\_medium} and $3.7\times$ on
\variant{warehouse\_narrow}, each at $p = 0.008$ --- the floor of a five-seed
permutation test --- and up to $4.6\times$ with entry admission also enabled.
On \variant{warehouse\_bottleneck} the two arms are identical, because no aisle
there ever commits a direction and the treatment has nothing to change.

It is the largest single effect measured anywhere in this system. \Cref{fig:forest} plots it beside every other
mechanism, and \Cref{sec:results-mechanisms} gives the table.

\begin{caveat}
The general form of this finding is worth stating separately from TOLL-PIBT,
because it applies to any local, push-based planner: \emph{a traffic rule
imposed as a hard constraint removes the freedom that local conflict resolution
runs on.} If the planner resolves blockages by displacement, every constraint
that shrinks the candidate set shrinks the set of blockages it can resolve.
Price the rule instead, above the preferences and below the objective.
\end{caveat}

\subsection{The aisle as a traffic signal: a maximum green, not only a minimum}
\label{sec:contribution-maxgreen}

An aisle decides its direction from the aggregate demand of the robots routed
through it, $S_a^{+}$ forward and $S_a^{-}$ reverse
(\Cref{sec:aisle-demand}), and commits when the imbalance leaves a dead band:
\begin{equation}
  \text{commit \textsc{forward} if } S_a^{+} - S_a^{-} > \tau_{\mathrm{switch}},
  \qquad
  \text{\textsc{reverse} if } S_a^{+} - S_a^{-} < -\tau_{\mathrm{switch}},
\end{equation}
with a minimum lock $T_{\min}$ after each commit. Hysteresis of this kind is
standard and it is only half a traffic signal: the dead band and the minimum
lock bound how \emph{soon} a direction may change, and nothing bounds how long
it may persist.

That gap is not a corner case in a warehouse; it is the normal operating
condition. Put pickups down one side of the floor and deliveries down the
other --- the standard layout, and three of the four maps used here --- and
loaded robots want one direction while empty robots want the other, in
near-equal measure. The imbalance sits inside the dead band forever, and an
aisle that committed a direction once holds it for the rest of the run while
half its traffic waits.

TOLL-PIBT adds a maximum green $T_{\max}$: past $T_{\max}$ steps in one
direction, \emph{any} opposing demand at all forces the aisle to drain and
flip. This gives a bounded-wait property that the dead band alone does not.

\begin{proposition}[Bounded direction wait]
\label{prop:maxgreen}
If an aisle $a$ holds direction $\delta$ from time $t_0$ and there is opposing
demand at every step, then $a$ commits $\bar{\delta}$ no later than
$t_0 + T_{\max} + T_{\mathrm{drain}}$, where $T_{\mathrm{drain}} \leq
\param{max\_drain\_time}$ is bounded because a \textsc{draining} aisle that has
not emptied within that many steps is force-reopened
(\Cref{sec:aisle-state-machine}).
\end{proposition}

Flips forced this way are counted separately as \code{starvation\_flips},
which keeps two questions apart that a single switch counter conflates: ``does
hysteresis stop oscillation?'' and ``does anybody starve?''

Measured, the maximum green is \emph{not} significant on throughput --- and
that is itself the finding. It was decisive while direction was a hard
constraint, because then a starved robot had no way out at all. Once
counterflow is merely expensive, a robot can buy its way out of a starved
aisle, so liveness no longer rests on the signal alone: the two fixes are
partly redundant. It stays enabled because it makes the aisle layer
\emph{correct} rather than merely survivable, and because its effect on
switching is significant: on \variant{warehouse\_narrow}, $22.5$ direction
switches per 1000 steps against $1.0$ without it ($p = 0.016$), of which $5.2$
per run are flips the rule forced rather than demand. Without it, those aisles
never turn at all. The third comparison animation shows
an aisle that never flips beside one that must.

\subsection{Recovery on corroborated stalls only}
\label{sec:contribution-corroboration}

A deadlock detector in a lifelong warehouse faces an unusual difficulty: in
saturation, the observable signature of a deadlock is also the observable
signature of a perfectly healthy queue. Three stall signals are available
(\Cref{sec:deadlock-signals}):
\begin{enumerate}
\item no route progress for $t_{\mathrm{blocked}}$ steps;
\item a cycle in the wait-for graph, where $i \to j$ when $i$'s
  priority-inheritance request went to $j$;
\item a repeated configuration --- the robot's last $k$ positions equal the $k$
  before them.
\end{enumerate}
Signal 1 alone describes every robot in every queue on a busy floor. Treating
it as sufficient means recovery fires constantly on healthy traffic, and
recovery's upper levels are drastic: temporary reverse movement, escape-vertex
assignment, waypoint hijack (\Cref{sec:recovery-levels}). Queues get taken
apart and rebuilt continuously.

TOLL-PIBT requires corroboration: a group must show signal 1 \emph{and} at least
one of signals 2 or 3 before escalation begins. Signals 2 and 3 are the ones
with a genuine deadlock semantics --- a circular wait, or a state the system
has demonstrably returned to. On \variant{warehouse\_corridors} this is worth
$0.134$ against $0.022$ tasks per timestep ($p = 0.008$), a six-fold difference
obtained entirely by declining to act on the weakest evidence. On the other
maps it points the same way without clearing five-seed significance. The fourth comparison
animation shows both.

There is a second-order lesson in how this defect stayed hidden. While
direction was a hard constraint, the architecture was manufacturing real
deadlocks continuously, so an over-eager detector looked useful: it \emph{was}
rescuing real deadlocks, of the system's own making. Fixing
\Cref{sec:contribution-soft} is what made the over-eagerness visible.

\subsection{Cost}
\label{sec:contribution-cost}

The aisle layer is a constant-factor addition to PIBT, not a change in its
complexity class. Per timestep, for fleet $A$ on graph $G$ with aisle set
$\mathcal{A}$:

\begin{table}[H]
\centering
\small
\begin{tabular}{@{}llr@{}}
\toprule
\textbf{Planner} & \textbf{Per-timestep cost} & \textbf{Measured (ms/step)} \\
\midrule
PIBT \citep{okumura2019pibt}
  & $O(|A|(\Delta(G) + F + \log|A|))$ & 0.4--1.0 \\
TOLL-PIBT (this work)
  & $+\,O(|A| + |\mathcal{A}|)$ aisle bookkeeping & 1.4--4.4 \\
Token Passing \citep{ma2017tp}
  & $O(|A| \cdot \mathrm{A}^\ast(|V| \cdot T_{\mathrm{horizon}}))$ & 32--222 \\
RHCR \citep{li2021rhcr}
  & windowed CBS, exponential in conflicts, capped & 43--98 \\
\bottomrule
\end{tabular}
\caption{Per-timestep cost. Measured columns are means over five seeds on the
three baseline maps of \Cref{sec:results}; the spread is across maps and fleet
sizes. The aisle layer costs a small constant factor over PIBT; both
path-planning baselines cost one to two orders of magnitude more.}
\label{tab:complexity}
\end{table}

The added work per timestep is: updating each aisle's demand sums
(one pass over robots), running the direction decision for each aisle (constant
work each), granting reservations in priority order (one pass), and the
deadlock detector's wait-for cycle search (linear in the fleet). None of it
involves a search, and none of it is repeated per candidate --- the congestion
and distance terms consumed by the score are precomputed or cached
(\Cref{sec:congestion,sec:graph-distances}), which is what keeps $F$ at $O(1)$.

\subsection{What this does not claim}
\label{sec:contribution-scope}

The honest summary, stated here rather than buried in \Cref{sec:results}: the
aisle layer is ahead where aisles are scarce and contended, and behind where
they are not. Against plain lifelong PIBT it gains 27\% on
\variant{warehouse\_bottleneck} and 21\% on \variant{warehouse\_corridors}, and
loses 18\% and 19\% on the two open maps --- and at five seeds \emph{the two
losses are significant and the two wins are not} ($p = 0.056$ and $0.063$
against $0.016$ and $0.008$). The fifth comparison animation is the losing
case, shown at the same size as the winning ones.

So the claim this document defends is precise: the mechanisms of this section
are established, the scoping of the aisle layer to aisle-shaped maps is
consistent and not yet significant, and \Cref{tab:honest} says which is which.
